%================================================================
% TEMPLATE PROPOSAL INOVASI HACKATHON
% IPB University
%
% Diadaptasi dari Template PPKI IPB University Edisi ke-4
% Disesuaikan khusus untuk format Proposal Kompetisi / Hackathon
% Batasan Halaman: 4--7 Halaman (termasuk Cover & Daftar Pustaka)
%
% Compiler yang direkomendasikan: pdfLaTeX
%================================================================

\documentclass[10pt,a4paper,oneside]{article}

%--------------------------------------------------------------
% # 1. Encoding & Bahasa
%--------------------------------------------------------------
\usepackage[T1]{fontenc}
\usepackage{textcomp}
\usepackage[utf8]{inputenc}
\usepackage[indonesian, provide=*]{babel}

%--------------------------------------------------------------
% # 2. Tipografi & Font (Times New Roman Standar Publikasi)
%--------------------------------------------------------------
\usepackage{amsmath}
\usepackage{amssymb}
\usepackage{mathptmx}
\usepackage{microtype}

%--------------------------------------------------------------
% # 3. Ukuran Halaman & Margin
% - Margin 2 cm pada semua sisi
% - Pemisah Kolom Jurnal (columnsep) 0.65 cm
%--------------------------------------------------------------
\usepackage[
  a4paper,
  top=2cm,
  left=2cm,
  right=2cm,
  bottom=2cm,
  headheight=15pt
]{geometry}
\setlength{\columnsep}{0.65cm}

% Parameter Float Jurnal Kompak (Cegah Halaman Kosong Akibat Float)
\renewcommand{\topfraction}{0.95}
\renewcommand{\bottomfraction}{0.85}
\renewcommand{\textfraction}{0.05}
\renewcommand{\floatpagefraction}{0.85}
\renewcommand{\dbltopfraction}{0.95}
\renewcommand{\dblfloatpagefraction}{0.85}

%--------------------------------------------------------------
% # 4. Spasi & Paragraf Gaya Jurnal
%--------------------------------------------------------------
\usepackage{setspace}
\setstretch{1.05}                % Spasi 1.05 standar publikasi ilmiah
\usepackage{indentfirst}         % Indentasi paragraf pertama
\setlength{\parindent}{0.4cm}    % Indentasi standar jurnal
\setlength{\parskip}{0pt}        % Tanpa jeda kosong antar paragraf gaya jurnal

%--------------------------------------------------------------
% # 5. Header & Footer (fancyhdr)
%--------------------------------------------------------------
\usepackage{fancyhdr}
\pagestyle{fancy}
\fancyhf{}
\fancyhead[L]{\fontsize{8pt}{10pt}\selectfont\textit{Proposal Inovasi TropiCare -- Hackathon 2026}}
\fancyhead[R]{\fontsize{9pt}{11pt}\selectfont\thepage}
\renewcommand{\headrulewidth}{0.4pt}
\renewcommand{\footrulewidth}{0pt}

% Halaman pertama (cover) tanpa header
\fancypagestyle{plain}{%
  \fancyhf{}
  \renewcommand{\headrulewidth}{0pt}
}

%--------------------------------------------------------------
% # 6. Format Judul Bab & Subbab (Gaya Jurnal Ilmiah)
%--------------------------------------------------------------
\usepackage{titlesec}

\titleformat{\section}
  {\fontsize{10.5pt}{12.5pt}\selectfont\bfseries}
  {\thesection}
  {0.6em}
  {\MakeUppercase}

\titleformat{\subsection}
  {\fontsize{10pt}{12pt}\selectfont\bfseries}
  {\thesubsection}
  {0.6em}
  {}

\titleformat{\subsubsection}
  {\fontsize{9.5pt}{11.5pt}\selectfont\bfseries\itshape}
  {\thesubsubsection}
  {0.6em}
  {}

\titlespacing*{\section}{0pt}{8pt}{3pt}
\titlespacing*{\subsection}{0pt}{6pt}{2pt}
\titlespacing*{\subsubsection}{0pt}{4pt}{1pt}

% Spasi Formula Matematika Kompak
\setlength{\abovedisplayskip}{3pt plus 1pt minus 1pt}
\setlength{\belowdisplayskip}{3pt plus 1pt minus 1pt}
\setlength{\abovedisplayshortskip}{1pt}
\setlength{\belowdisplayshortskip}{1pt}

% Format penomoran: BAB I -> 1.1 -> 1.1.1
\renewcommand{\thesection}{BAB \Roman{section}}
\renewcommand{\thesubsection}{\arabic{section}.\arabic{subsection}}
\renewcommand{\thesubsubsection}{\arabic{section}.\arabic{subsection}.\arabic{subsubsection}}

%--------------------------------------------------------------
% # 7. Tabel & Gambar
%--------------------------------------------------------------
\usepackage{booktabs}        % \toprule, \midrule, \bottomrule
\usepackage{array}
\usepackage{tabularx}
\usepackage{multirow}
\renewcommand{\arraystretch}{0.92} % Kerapatan baris tabel
\usepackage{graphicx}
\usepackage{float}
\usepackage{tikz}
\usetikzlibrary{positioning, arrows.meta, shapes.geometric, backgrounds, fit}

\usepackage{caption}
\DeclareCaptionLabelFormat{ipblabel}{#1~#2}
\captionsetup{
  format          = hang,
  labelformat     = ipblabel,
  labelsep        = space,
  font            = small,
  labelfont       = bf,
  textfont        = {},
  justification   = justified,
  singlelinecheck = true
}
\captionsetup[table]{name=Tabel, position=above, skip=3pt}
\captionsetup[figure]{name=Gambar, position=below, skip=4pt}

%--------------------------------------------------------------
% # 8. Kotak Sorotan / Callout Box Modern (tcolorbox)
%--------------------------------------------------------------
\usepackage[most]{tcolorbox}
\definecolor{ipbblue}{RGB}{20, 52, 114}
\definecolor{accentteal}{RGB}{0, 128, 128}
\definecolor{lightbg}{RGB}{248, 250, 252}

\newtcolorbox{highlightbox}[1][]{
  colback=lightbg,
  colframe=ipbblue,
  fonttitle=\bfseries\small,
  coltitle=white,
  boxrule=0.6pt,
  arc=1.5mm,
  left=6pt,
  right=6pt,
  top=4pt,
  bottom=4pt,
  #1
}

%--------------------------------------------------------------
% # 9. Referensi & Sitasi (Gaya IEEE Numerik [1], [2])
%--------------------------------------------------------------
\usepackage[numbers,sort&compress]{natbib}
\setlength{\bibsep}{1.2pt}

%--------------------------------------------------------------
% # 10. Hyperlink
%--------------------------------------------------------------
\usepackage[
  colorlinks=true,
  linkcolor=ipbblue,
  citecolor=ipbblue,
  urlcolor=ipbblue,
  pdfauthor={Tim Pengusul IPB},
  pdftitle={Proposal Inovasi Hackathon}
]{hyperref}

% Kontrol List / Poin
\usepackage{enumitem}
\setlist{nosep, leftmargin=1.2em}

%==============================================================
% # DOKUMEN UTAMA
%==============================================================
\begin{document}

% 1. Halaman Sampul (Cover)
%==============================================================
% HALAMAN JUDUL (COVER) PROPOSAL
% TropiCare - Hackathon JOINTSXINSPIRE 2026 (UGM)
%==============================================================

\begin{titlepage}
\thispagestyle{empty}
\begin{center}

\vspace*{0.5cm}

% Identitas Kompetisi
{\fontsize{11pt}{14pt}\selectfont\bfseries
PROPOSAL INOVASI TEKNOLOGI KESEHATAN\\
HACKATHON JOINTSXINSPIRE 2026\\
UNIVERSITAS GADJAH MADA\par
}

\vspace{1.2cm}

% Judul Utama
{\fontsize{24pt}{28pt}\selectfont\bfseries TropiCare\par}

\vspace{0.5cm}

{\fontsize{11.5pt}{15pt}\selectfont\bfseries
Sistem Triase Mandiri Pra-Hospital Berbasis Retrieval-Augmented Large Language Model untuk Skrining Komunitas dan Navigasi Faskes Layanan Primer\par
}

\vspace{0.6cm}

{\fontsize{10pt}{13pt}\selectfont\itshape
Penapisan komunitas berbasis privasi Zero-NIK, direktori nasional 13.930 faskes SATUSEHAT MSI, dan rekayasa perangkat lunak berdisiplin tinggi.\par
}

\vspace{1.4cm}

% Logo IPB University
\includegraphics[width=3cm]{images/ipb_logo}

\vspace{1.4cm}

% Tim Pengusul
{\fontsize{10.5pt}{14pt}\selectfont
\textbf{Diusulkan oleh:}\\[0.4cm]
\begin{tabular}{ll}
\textbf{Walid Nadirul Ahnaf} & -- \textit{Ilmu Komputer, SSMI IPB}\\
\textbf{Petrus Yoga Haleserens} & -- \textit{Kedokteran, FK IPB}\\
\textbf{Aldi Pramudya} & -- \textit{Ilmu Komputer, SSMI IPB}\\
\textbf{Erlangga Sri Heryanto} & -- \textit{Statistika, SSMI IPB}\\
\end{tabular}\par
}

\vfill

{\fontsize{11pt}{15pt}\selectfont\bfseries
INSTITUT PERTANIAN BOGOR\\
BOGOR\\
2026\par
}

\end{center}
\end{titlepage}
\newpage


\clearpage
\twocolumn

% Mulai penomoran halaman isi
\setcounter{page}{2}

% 2. Bab I: Pendahuluan
%==============================================================
% BAB I PENDAHULUAN
%==============================================================

\section{Pendahuluan}

\subsection{Latar Belakang dan Urgensi Masalah}

[Tuliskan latar belakang permasalahan di sini. Jelaskan konteks kesehatan di Indonesia, beban ganda penyakit tropis dan penyakit tidak menular, serta hambatan akses di fasilitas kesehatan tingkat pertama (FKTP) dan wilayah kepulauan/pedalaman. Cantumkan data pendukung dan referensi seperti \citep{Kemenkes2024}].

\begin{enumerate}
    \item \textbf{[Masalah 1 - Antrean dan Administrasi FKTP]:} [Uraikan kendala antrean panjang di Puskesmas, kasus ringan SKDI 4A yang seharusnya bisa mandiri, dan pembatalan tes lab karena pasien tidak puasa].
    \item \textbf{[Masalah 2 - Akses Geografis dan Salah Rujukan]:} [Uraikan kendala transportasi wilayah kepulauan/pedalaman dan risiko salah rujukan ke faskes tanpa stok obat program].
    \item \textbf{[Masalah 3 - Rendahnya Skrining Mandiri Usia Produktif]:} [Uraikan rendahnya pemeriksaan berkala tensi, gula darah, dan antropometri di masyarakat tanpa alat medis].
    \item \textbf{[Masalah 4 - Bahaya Swamedikasi dan Tanggap Darurat]:} [Uraikan risiko konsumsi obat bebas keliru seperti NSAID pada demam dengue dan keterlambatan akses darurat].
\end{enumerate}

\subsection{Rumusan Masalah}

[Rumuskan pertanyaan kunci yang dijawab oleh inovasi ini]:
\begin{enumerate}
    \item Bagaimana mengembangkan sistem triase mandiri berbasis *Medical Large Language Model* yang terikat bukti klinis (*evidence-grounded*) untuk memilah keparahan kondisi tanpa risiko halusinasi?
    \item Bagaimana merancang integrasi direktori fasilitas kesehatan nasional SATUSEHAT MSI dan rujukan terdekat berbasis privasi tanpa meminta data kependudukan sensitif?
    \item Bagaimana menyediakan deteksi interaksi obat aktif dan proksi antropometri fisik bagi masyarakat yang tidak memiliki instrumen medis rumah tangga?
\end{enumerate}

\subsection{Tujuan Inovasi}

[Tuliskan tujuan spesifik dan terukur dari inovasi yang diusulkan]:
\begin{enumerate}
    \item Mewujudkan sistem triase mandiri pra-hospital 4 rute (Merah, Kuning Khusus, Kuning, Hijau) berstandar PPK Dokter di FKTP Kemenkes RI.
    \item Mengotomasi navigasi rujukan faskes terdekat ber-WhatsApp dan transmisi ringkasan klinis terenkripsi (*QR Code*).
    \item Memberikan edukasi keamanan swamedikasi obat dan kepatuhan persiapan pra-laboratorium bagi pasien.
\end{enumerate}


% 3. Bab II: Solusi dan Inovasi yang Ditawarkan
%==============================================================
% BAB II SOLUSI DAN INOVASI YANG DITAWARKAN
%==============================================================

\section{Solusi dan Inovasi yang Ditawarkan}

\subsection{Deskripsi Solusi dan Alur Klinis Berjenjang}

TropiCare adalah aplikasi web skrining awal penyakit tropis dan penyakit jangka panjang (PTM) yang mudah diakses tanpa instalasi dan tanpa kewajiban membuat akun. Seluruh logika medis berpedoman pada Keputusan Menteri Kesehatan RI No. HK.01.07/MENKES/1186/2022 tentang Panduan Praktik Klinis bagi Dokter di FKTP \citep{Kemenkes2022a} serta Permenkes No. 3 Tahun 2026 tentang Penanggulangan Penyakit \citep{Permenkes2026}. Sistem bekerja dalam waktu sekitar 3 menit melalui empat lapisan penapisan:

\begin{enumerate}
    \item \textbf{Gerbang Privasi, Data Diri, \& Deteksi Cepat 11 Tanda Bahaya (\textit{Red Flags}):} 
    Setelah memilih jalur skrining (infeksi akut atau risiko jangka panjang) dan mencentang persetujuan privasi, pengguna mengisi data diri dasar (usia, jenis kelamin, dan status kehamilan) lalu sistem langsung memeriksa 11 kondisi gawat darurat: (1) muntah hebat atau tangan-kaki dingin (syok dengue), (2) perdarahan sulit berhenti atau batuk darah banyak, (3) tidak buang air kecil $>6\text{--}8$ jam atau urin gelap (gagal ginjal akut), (4) sesak napas sangat berat, (5) penurunan kesadaran atau kejang, (6) sangat lemas dan mata cekung (dehidrasi berat), (7) luka gigitan hewan penular rabies, (8) nyeri dada kiri menjalar (serangan jantung), (9) tanda stroke mendadak (SeGeRa Ke RS), (10) krisis darah tinggi, serta (11) tanda koma diabetes \citep{Kemenkes2022a}. Satu saja tanda bahaya terpilih (atau muncul mata/kulit sangat kuning saat demam pada gejala penyerta), kuesioner langsung berhenti dan layar berubah menjadi Rute Merah tenang tanpa animasi yang menampilkan tombol telepon 119, nomor darurat 119/112 yang dapat disalin, serta panduan pertolongan pertama (\textit{DO} dan \textit{DON'T}) sambil menuju rumah sakit. Deteksi ini dieksekusi di peramban ponsel dalam $<500\text{ ms}$ walau peladen sedang mati.

    \item \textbf{Kuesioner Adaptif, Peta Ekologi (\textit{One Health}), \& Peta Nyeri:} 
    Jika tidak ada tanda bahaya, pertanyaan menyesuaikan jawaban sebelumnya ke dalam 5 klaster infeksi tropis (demam, napas, pencernaan, kulit, bengkak) untuk 20 penyakit tropis atau 3 klaster penyakit kronis (gula \& saraf, tekanan \& kepala, sendi) untuk 6 kondisi PTM \citep{PERKENI2024}. Sistem menggali riwayat lingkungan 1--4 bulan terakhir (daerah endemis malaria di NTT/Papua atau skistosomiasis di Sulawesi Tengah, genangan banjir leptospirosis, kontak TBC, lingkungan padat skabies) \citep{Permenkes2026}, dilengkapi peta nyeri tubuh interaktif dan skala nyeri 0--10 berbasis wajah.

    \item \textbf{Pemeriksaan Keamanan Obat \& Proksi Fisik Tanpa Alat Medis:} 
    Saat pengguna mengetik atau memilih obat yang sudah diminum, sistem memeriksa keamanan obat (4 aturan intervensi aktif ditambah 3 skema kewaspadaan komorbid dengue). Suspek demam berdarah yang meminum NSAID (Ibuprofen, Aspirin, Asam Mefenamat) langsung mendapat peringatan kritis merah untuk menghentikannya demi mencegah perdarahan lambung \citep{Alam2022,Mohiuddin2023,Tovar2024}; penggunaan obat penahan diare (Loperamide) pada diare berdarah memicu peringatan kuning \citep{Koo2007,Schiller2020}; sedangkan bengkak kaki akibat Amlodipine dan batuk kering akibat Captopril diberi catatan informasi biru agar tidak disalahartikan sebagai filariasis atau TBC \citep{Liang2022,Savage2020,Sica2003,Dicpinigaitis2014}. Bagi rumah tangga tanpa termometer atau timbangan \citep{Kemenkes2024}, tersedia pertanyaan pengganti (menggigil hebat, nadi cepat saat istirahat, dan perubahan kelonggaran pakaian) serta input tanda vital yang otomatis menaikkan rute ke Kuning bila suhu $\ge 39,5^\circ\text{C}$ atau tensi sistolik $>160\text{ mmHg}$ \citep{PERKENI2024}.

    \item \textbf{Orkestrasi 4 Rute Layanan \& Tindak Lanjut Faskes:} 
    Hasil skrining menyajikan satu perkiraan utama beserta skor kecocokan (atau \textit{``belum dapat ditentukan''} yang diarahkan ke Rute Kuning bila skor $<30\%$), tiga kemungkinan lain, alasan \textit{``Kenapa aku menebak ini?''}, serta satu dari empat rute layanan:
    \begin{itemize}
        \item \textbf{Rute Merah (IGD Rumah Sakit \& 119):} Kondisi darurat diarahkan langsung ke IGD RS terdekat atau Rabies Center.
        \item \textbf{Rute Kuning Khusus (Poli Program Puskesmas):} Dugaan penyakit program nasional (TBC, Kusta, Frambusia, Filariasis, Skistosomiasis, atau Campak) diarahkan ke Puskesmas penyelenggara program.
        \item \textbf{Rute Kuning (Puskesmas / FKTP):} Perlu pemeriksaan dokter dan laboratorium; pada jalur kronis dilengkapi instruksi wajib puasa 10--12 jam (hanya minum air putih) sebelum cek gula darah dan kolesterol \citep{PERKENI2024}.
        \item \textbf{Rute Hijau (Rawat Mandiri):} Gejala ringan non-komplikasi dipandu perawatan mandiri di rumah dan apotek terdekat.
    \end{itemize}
\end{enumerate}

\subsection{Arsitektur Sistem dan Alur Data Transparan}

Alur kerja sistem dari perangkat pasien hingga layanan faskes dirancang secara transparan dengan memisahkan pemrosesan di perangkat, di peladen TropiCare, dan di layanan AI eksternal (Gambar~\ref{fig:arsitektur}).

\begin{figure}[H]
\centering
\resizebox{0.96\textwidth}{!}{%
\begin{tikzpicture}[
    node distance=0.55cm and 0.45cm,
    box/.style={rectangle, draw=black!70, fill=black!4, rounded corners=2mm, text centered, font=\footnotesize, minimum height=0.85cm, line width=0.7pt},
    cloud/.style={rectangle, draw=black, fill=black!10, rounded corners=2mm, text centered, font=\footnotesize\bfseries, minimum height=0.85cm, line width=0.8pt},
    arrow/.style={-{Stealth[scale=0.8]}, draw=black!70, line width=0.7pt}
]
    % Top pipeline
    \node[box, text width=4.3cm] (input) {\textbf{Di Perangkat Pasien (React)}\\Persetujuan Privasi, Lokasi $\rightarrow$ Provinsi, \& Cek 11 Tanda Bahaya};
    \node[cloud, text width=4.6cm, right=of input] (engine) {\textbf{Peladen TropiCare (FastAPI)}\\Endpoint Stateless, Aturan Rute, \& Mirror 13.930 Faskes};
    \node[box, text width=4.3cm, right=of engine] (output) {\textbf{Layar Hasil \& Tindak Lanjut}\\4 Rute, 5 Faskes Terdekat, Mobile JKN, \& QR Petugas};

    % Bottom pipeline
    \node[box, text width=4.3cm, below=0.7cm of input] (red) {\textbf{Layar MERAH + 119}\\Jalan langsung di browser walau peladen mati};
    \node[box, text width=4.6cm, below=0.7cm of engine] (ai) {\textbf{Layanan AI (Gemma 4 / MedGemma)}\\Skor kecocokan + RAG 203 kutipan \& kenali teks obat};
    \node[box, text width=4.3cm, below=0.7cm of output] (qr) {\textbf{PostgreSQL \& Pemindai QR}\\Statistik tanpa identitas langsung; rangkuman dibuka petugas};

    % Arrows
    \draw[arrow] (input) -- (engine);
    \draw[arrow] (engine) -- (output);
    \draw[arrow] (input) -- (red);
    \draw[arrow] (engine) -- (ai);
    \draw[arrow] (ai) -- (engine);
    \draw[arrow] (output) -- (qr);
\end{tikzpicture}%
}
\caption{Alur Sistem TropiCare dari Perangkat Pasien, Peladen FastAPI, hingga Layanan AI dan Petugas Faskes}
\label{fig:arsitektur}
\end{figure}

Keunggulan rekayasa produk TropiCare bertumpu pada lima pilar teknis:
\begin{enumerate}
    \item \textbf{Evaluasi Darurat di Peramban (\textit{Client-Side Safety Gate}):} Pemeriksaan 11 tanda bahaya darurat dan pembuatan kode QR berjalan langsung di peramban ponsel. Jika tanda bahaya terpilih, layar darurat Merah dan tombol 119 aktif seketika meskipun koneksi internet atau peladen sedang mati.
    \item \textbf{Privasi Terjaga (\textit{Zero-NIK} \& Lokasi di Perangkat):} Pengguna dapat menyelesaikan skrining sebagai tamu tanpa akun dan tanpa pernah dimintai NIK atau nama. Koordinat GPS hanya diolah di memori peramban bersama poligon batas wilayah untuk mengenali provinsi dan mengurutkan jarak 5 faskes terdekat; hanya nama provinsi yang dikirim ke peladen \citep{UUPDP2022}.
    \item \textbf{Mirror Lokal 13.930 Faskes \& Jembatan Mobile JKN:} Peladen memuat mirror lokal 13.930 faskes resmi SATUSEHAT (10.374 Puskesmas dan 3.556 Rumah Sakit). Daftar 5 faskes terdekat dilengkapi tautan Google Maps, tombol kirim rangkuman via WhatsApp (atau telepon), serta tombol pembuka aplikasi Mobile JKN untuk pendaftaran janji dokter.
    \item \textbf{Penalaran Gemma 4 / MedGemma + RAG yang Terikat Pedoman Dokter:} Jawaban dianalisis menggunakan model keluarga Gemma (\textbf{Gemma 4 / MedGemma}) atau penyedia kompatibel (OpenRouter, vLLM, maupun Gemini) yang diikat pada 203 kutipan literatur medis terverifikasi (KMK 1186/2022 dan pedoman penyakit tropis Kemenkes) melalui indeks leksikal BM25 pra-komputasi. AI juga membantu mengenali nama obat ketikan pasien, sedangkan keputusan rute darurat dan ambang vital tetap dikunci oleh aturan deterministik.
    \item \textbf{QR Rangkuman untuk Petugas Faskes:} Jawaban pasien dikodekan secara ringkas ke dalam kode QR di layar ponsel tanpa menyimpan identitas di peladen. Saat dipindai petugas Puskesmas/klinik, rangkuman kondisi, riwayat obat, dan saran pemeriksaan laboratorium langsung terbuka di peramban petugas tanpa perlu wawancara ulang dari nol.
\end{enumerate}

\subsection{Perbandingan dengan Alternatif yang Ada}

TropiCare dirancang untuk \textbf{melengkapi, bukan menggantikan} layanan resmi seperti Mobile JKN: setelah tahu harus ke mana dari hasil skrining TropiCare, pasien langsung menekan tombol Mobile JKN untuk mendaftar antrean dokter. Perbandingan fitur disajikan pada Tabel~\ref{tab:keunggulan}.

\begin{table}[H]
\centering
\footnotesize
\caption{Perbandingan TropiCare dengan Alternatif Aplikasi Kesehatan di Indonesia}\label{tab:keunggulan}
\vspace{3pt}
\begin{tabularx}{\textwidth}{p{3.6cm} X X X}
\toprule
\textbf{Aspek Pembanding} & \textbf{TropiCare} & \textbf{Mobile JKN (BPJS)} & \textbf{Ada Health} \\
\midrule
\textbf{Skrining Gejala Infeksi Tropis} & \textbf{Ya (20 penyakit tropis Indonesia)} & Tidak tersedia & Umum, tidak khusus penyakit tropis Indonesia \\
\midrule
\textbf{Skrining Risiko Kronis (PTM)} & \textbf{Ya (6 kondisi PTM + proksi fisik)} & Ya (47 pertanyaan tetap, dibatasi $1\times$/tahun) & Umum \\
\midrule
\textbf{Kuesioner Adaptif} & \textbf{Ya (menyesuaikan jawaban sebelumnya)} & Tidak (daftar pertanyaan tetap) & Ya \\
\midrule
\textbf{Tanda Bahaya $\rightarrow$ Darurat Langsung} & \textbf{Ya (jalan di browser walau server mati)} & Tidak tersedia & Saran umum menghubungi darurat \\
\midrule
\textbf{Arahan Jenis Layanan} & \textbf{4 Rute (IGD, Poli Program, Puskesmas, Mandiri)} & Tingkat risiko \& saran gaya hidup & Saran langkah berikutnya secara umum \\
\midrule
\textbf{Rangkuman Petugas (QR) \& Bahasa} & \textbf{QR rangkuman instan; Bahasa Indonesia} & Tanpa QR pra-triase; Bahasa Indonesia & Laporan asesmen; belum tersedia Bahasa Indonesia konsumen \\
\bottomrule
\end{tabularx}
\end{table}


% 4. Bab III: Target Pengguna dan Manfaat yang Diharapkan
%==============================================================
% BAB III TARGET PENGGUNA DAN MANFAAT YANG DIHARAPKAN
%==============================================================

\section{Target Pengguna dan Manfaat yang Diharapkan}

\subsection{Target Pengguna Solusi}

TropiCare melayani empat kelompok pengguna dalam alur layanan kesehatan primer:

\begin{enumerate}
    \item \textbf{Masyarakat Umum dan Kelompok Usia Produktif (15--64 Tahun):} 
    Warga yang mengalami gejala infeksi akut atau memiliki faktor risiko penyakit tidak menular (hipertensi, hiperglikemia, obesitas sentral). Kelompok ini sering menunda pemeriksaan karena jam kerja dan antrean panjang FKTP. Lewat TropiCare, pengguna dapat melakukan skrining mandiri melalui peramban ponsel tanpa mengunduh aplikasi dan tanpa membuat akun, serta menjawab pertanyaan proksi fisik (pola menggigil, nadi istirahat, dan perubahan ukuran pakaian) bila tidak memiliki alat ukur medis di rumah \citep{Kemenkes2024}.
    
    \item \textbf{Keluarga Pendamping Pasien Komorbid, Lansia, dan Anak:} 
    Anggota keluarga atau pengasuh (\textit{caregiver}) yang merawat lansia dengan penyakit kronis maupun anak yang sedang demam atau diare. Sistem memeriksa obat yang diminum pasien untuk mencegah bahaya salah obat bebas (seperti larangan NSAID pada anak suspek demam berdarah dan larangan obat penahan diare pada disentri \citep{Alam2022,Koo2007}), sekaligus mengingatkan syarat puasa 10--12 jam sebelum mengantar lansia cek darah di Puskesmas \citep{PERKENI2024}.
    
    \item \textbf{Masyarakat Wilayah Kepulauan, Pedalaman, dan Daerah 3T:} 
    Warga di daerah endemis penyakit tropis terabaikan, seperti Lembah Napu dan Lindu di Sulawesi Tengah (Skistosomiasis) serta NTT dan Papua (Malaria dan Frambusia). Rute Kuning Khusus mengarahkan warga langsung ke Puskesmas penyelenggara program eliminasi beserta kontak WhatsApp/teleponnya, sehingga warga tidak membuang biaya perjalanan ke faskes yang tidak menyimpan stok obat program.
    
    \item \textbf{Tenaga Kesehatan dan Pengelola FKTP (Puskesmas dan Klinik Pratama):} 
    Dokter, perawat, petugas loket pendaftaran, dan kader Posyandu/Posbindu. Petugas memindai kode QR pasien untuk membuka rangkuman keluhan, riwayat obat, dan saran pemeriksaan laboratorium di peramban tanpa perlu mewawancarai dan mengetik ulang keluhan dari awal.
\end{enumerate}

\subsection{Manfaat yang Diharapkan}

Penggunaan TropiCare memberikan tiga manfaat terukur bagi pasien, fasilitas kesehatan, dan pengelola program kesehatan:

\begin{enumerate}
    \item \textbf{Manfaat bagi Pasien dan Keluarga:}
    Pasien memperoleh panduan arah layanan dalam waktu sekitar 3 menit tanpa biaya; mengenali 11 tanda bahaya darurat dalam $<500\text{ ms}$ untuk segera menghubungi 119 atau menuju IGD; menghindari perdarahan lambung akibat NSAID saat demam berdarah serta \textit{toxic megacolon} akibat Loperamide saat disentri; dan datang ke laboratorium FKTP dalam kondisi sudah berpuasa 10--12 jam.

    \item \textbf{Manfaat bagi Tenaga Kesehatan dan Puskesmas:}
    Puskesmas dapat mengurangi tumpukan antrean sebesar 30\%--40\% karena pasien bergejala ringan non-komplikasi (SKDI 4A) teredukasi untuk menjalani perawatan mandiri (Rute Hijau) di rumah. Bagi pasien yang datang ke Puskesmas, pemindaian kode QR rangkuman mempercepat proses pra-triase di loket dan membantu dokter memilih pemeriksaan penunjang yang relevan.

    \item \textbf{Manfaat bagi Dinas Kesehatan dan Program Nasional:}
    Skrining awal faktor risiko PTM (hipertensi, diabetes, dislipidemia, dan risiko kardiovaskular) membantu menekan komplikasi pembiayaan katastropik BPJS Kesehatan. Di tingkat wilayah, Dinas Kesehatan dan Ditjen P2P Kemenkes RI memperoleh statistik agregat tren keluhan per provinsi tanpa identitas langsung (dengan penyembunyian kelompok berjumlah $<10$ untuk mencegah pelacakan balik individu) sesuai prinsip pelindungan data UU No. 27 Tahun 2022 \citep{UUPDP2022}.
\end{enumerate}


% 5. Bab IV: Rencana Implementasi dan Pengembangan
%==============================================================
% BAB IV SKALABILITAS, KELAYAKAN, DAN RENCANA PENGEMBANGAN
%==============================================================

\section{Skalabilitas, Kelayakan, dan Rencana Pengembangan}

\subsection{Arsitektur Skalabilitas Tiga Dimensi (Teknis, Konten, dan Jangkauan)}

Agar dapat tumbuh dari purwarupa fungsional menjadi layanan yang melayani ribuan pengguna di seluruh Indonesia tanpa membebani infrastruktur, TropiCare dirancang dengan skalabilitas pada tiga dimensi utama:

\begin{enumerate}
    \item \textbf{Skalabilitas Teknis (\textit{Stateless \& Modular}):} 
    Endpoint skoring pada backend FastAPI bersifat \textit{stateless} sehingga dapat direplikasi secara horizontal dengan mudah. Arsitektur menggunakan pola \textit{Ports and Adapters} (heksagonal): mesin skoring aturan deterministik maupun penyedia model AI (Gemma 4 / MedGemma via OpenRouter, vLLM mandiri, atau Gemini) dipasang pada lapisan komposisi (\textit{composition root}) dan dapat diganti hanya melalui konfigurasi tanpa mengubah logika bisnis. Selain itu, deteksi 11 tanda bahaya darurat dijalankan sepenuhnya di peramban pasien (biaya komputasi server nol), dan pratinjau skor selama pengisian kuis memakai mesin aturan lokal serta \textit{cache} memori 24 jam sehingga pemanggilan LLM hanya terjadi satu kali di akhir sesi skrining non-darurat.

    \item \textbf{Skalabilitas Konten Klinis (\textit{Data-Driven JSON}):} 
    Seluruh pertanyaan, opsi jawaban, aturan peringatan obat, dan bobot penyakit disimpan dalam berkas konfigurasi JSON terstruktur yang terpisah dari kode antarmuka React. Dokter dan statistisi dapat menambahkan penyakit baru, memperbarui bobot gejala, atau menyesuaikan profil endemisitas wilayah secara langsung tanpa perlu memodifikasi kode tampilan frontend. Setiap perubahan teks medis dan aturan klinis dikendalikan melalui persetujuan kolaborator dokter.

    \item \textbf{Skalabilitas Jangkauan \& Surveilans Berprivasi:} 
    Ekspansi layanan dirancang melalui lima tahap berurutan: \textbf{(1) Validasi Klinis} $\rightarrow$ \textbf{(2) Pilot Klinik \& Puskesmas} $\rightarrow$ \textbf{(3) TropiCare Plus (Pemantauan PTM)} $\rightarrow$ \textbf{(4) White-label untuk LSM Kesehatan} $\rightarrow$ \textbf{(5) API Mitra \& Surveilans Dinas Kesehatan}. Untuk mendukung Dinas Kesehatan tanpa melanggar privasi warga, statistik wilayah disajikan murni dari agregat tanpa identitas langsung dengan aturan privasi ketat: \textbf{kelompok agregat dengan jumlah kasus $<10$ disembunyikan} guna mencegah pelacakan balik individu.
\end{enumerate}

\subsection{Analisis Kelayakan Finansial (\textit{Feasibility}) dan Ekonomi Unit}

Keunggulan utama desain arsitektur ringan TropiCare adalah biaya operasional bulanan yang sangat rendah, sehingga layanan ini dapat bertahan hidup secara mandiri tanpa harus menunggu anggaran pengadaan pemerintah yang panjang. Rincian perhitungan biaya infrastruktur bulanan dan model keberlanjutan disajikan pada Tabel~\ref{tab:feasibility}.

\begin{table}[H]
\centering
\footnotesize
\caption{Rincian Biaya Operasional Bulanan dan Skema Monetisasi Berkelanjutan TropiCare}\label{tab:feasibility}
\vspace{3pt}
\begin{tabularx}{\textwidth}{p{3.4cm} p{2.8cm} X}
\toprule
\textbf{Komponen Biaya / Layanan} & \textbf{Nilai Finansial} & \textbf{Dasar Perhitungan Teknis \& Model Bisnis} \\
\midrule
\multicolumn{3}{l}{\textbf{A. Struktur Biaya Operasional Aplikasi (\textit{Monthly Burn Rate})}} \\
\midrule
\textbf{Server \& Basis Data\newline(\textit{Railway Hobby Plan})} & \textbf{Rp90.000 / bulan} & Menjalankan kontainer backend FastAPI \textit{stateless} dan basis data PostgreSQL (menyimpan 13.930 mirror faskes, riwayat akun opsional, dan agregat statistik). \\
\midrule
\textbf{Inferensi AI\newline(\textit{OpenRouter API})} & \textbf{Rp180.280 / bulan} & Melayani panggilan tunggal terstruktur (\textit{single-call JSON}) model keluarga Gemma + pengenalan obat. Dengan $\sim 1.550$ token/skrining dan \textit{cache} 24 jam, anggaran ini mencukupi $\sim \mathbf{15.000\text{--}20.000}$ skrining penuh/bulan ($\approx \text{Rp}9\text{--}\text{Rp}12$ per pasien). \\
\midrule
\textbf{Domain \& DNS Resmi} & \textbf{Rp25.000 / bulan} & Biaya tahunan nama domain resmi (\text{Rp}300.000/tahun) dibagi rata per bulan. \\
\midrule
\textbf{Total Biaya Server} & \textbf{Rp295.280 / bulan} & \textbf{Titik Impas (\textit{Break-Even}): Cukup 1 langganan klinik mitra (Rp400.000/bulan) sudah menutup seluruh biaya server sebulan penuh.} \\
\midrule
\multicolumn{3}{l}{\textbf{B. Skema Keberlanjutan Pendapatan (Skrining Inti Masyarakat Tetap Gratis)}} \\
\midrule
\textbf{1. Dasbor Faskes\newline(Klinik \& Puskesmas)} & \textbf{Rp250.000--Rp400.000\newline/ bulan / faskes} & Layanan B2B pra-triase loket lewat pemindaian QR dan rekapitulasi keluhan pasien datang untuk memangkas antrean administrasi. \\
\midrule
\textbf{2. Lisensi API \& Kemitraan Surveilans} & \textbf{Rp15.000.000--Rp40.000.000 / tahun} & Akses agregat tren penyakit per wilayah tanpa identitas langsung (kelompok $<10$ disembunyikan) bagi LSM/NGO eliminasi penyakit tropis, institusi riset, atau kemitraan CSR kesehatan. \\
\midrule
\textbf{3. TropiCare Plus\newline(Pantau PTM Mandiri)} & \textbf{Rp15.000 / bulan\newline(Opsional)} & Fitur tambahan bagi penderita hipertensi atau diabetes untuk memantau grafik tren tekanan darah dan gula darah jangka panjang. \\
\bottomrule
\end{tabularx}
\end{table}

\subsection{Peta Jalan Pengembangan, Mitigasi Risiko, dan Indikator Kinerja (KPI)}

Pelaksanaan pengembangan selama 12 bulan dari tahap validasi hingga integrasi ekosistem, beserta matriks mitigasi risiko dan indikator kinerja utama (KPI), dirangkum pada Tabel~\ref{tab:roadmap} dan Tabel~\ref{tab:kpi}.

\begin{table}[H]
\centering
\footnotesize
\caption{Peta Jalan Pengembangan 12 Bulan dan Matriks Mitigasi Risiko Utama}\label{tab:roadmap}
\vspace{3pt}
\begin{tabularx}{\textwidth}{p{2.5cm} p{3.1cm} X p{3.8cm}}
\toprule
\textbf{Tahap \& Periode} & \textbf{Fokus \& Risiko Utama} & \textbf{Aktivitas Eksekusi \& Mitigasi Konkret} & \textbf{Luaran Kunci (\textit{Milestone})} \\
\midrule
\textbf{Tahap 1\newline(Bulan 1--3)} & Validasi Klinis \& Mitigasi Risiko Halusinasi AI & Uji buta 100 kartu kasus dokter; kalibrasi bobot aturan JSON bersama statistisi; pemutusan 11 tanda bahaya langsung di browser tanpa AI (\textit{safety gate}); uji kemudahan pakai (SUS). & Akurasi triase $>92\%$ tanpa kasus darurat terlewat; skor SUS $>80$ (\textit{Grade A}); 304 kasus uji backend lulus. \\
\midrule
\textbf{Tahap 2\newline(Bulan 4--6)} & Pilot Klinik/FKTP \& Mitigasi Beban Loket & Pilot di 10 klinik pratama dan Puskesmas wilayah endemis; penerapan pemindaian QR rangkuman tanpa ketik ulang (\textit{zero-retyping}); kurasi nomor WhatsApp aktif faskes. & Operasional di 10 faskes pilot; biaya server tertutup penuh dari langganan mitra awal (\textit{self-sustaining}). \\
\midrule
\textbf{Tahap 3\newline(Bulan 7--9)} & TropiCare Plus, Regulasi \& Privasi Data & Peluncuran modul grafik Pantau PTM; pendaftaran \textit{Regulatory Sandbox} Kemenkes RI \citep{UUKesehatan2023,KemenkesDTO2024}; audit kepatuhan UU PDP No. 27/2022 (Zero-NIK, lokasi tetap di perangkat) \citep{UUPDP2022}. & Terdaftar di \textit{Regulatory Sandbox} Kemenkes; modul Pantau PTM aktif; jembatan rujukan Mobile JKN teruji. \\
\midrule
\textbf{Tahap 4--5\newline(Bulan 10--12)} & White-Label LSM \& API Mitra Surveilans & Kemitraan white-label dengan LSM pemberantasan penyakit tropis (TBC, Kusta, NTD); pembukaan API agregat wilayah bagi Dinkes (kelompok $<10$ disembunyikan); persiapan pemetaan standar FHIR R4. & Adopsi di 100+ faskes; dasbor surveilans agregat wilayah beroperasi; kesiapan interoperabilitas nasional. \\
\bottomrule
\end{tabularx}
\end{table}

\begin{table}[H]
\centering
\footnotesize
\caption{Indikator Kinerja Utama (KPI) Keberhasilan Implementasi TropiCare}\label{tab:kpi}
\vspace{3pt}
\begin{tabularx}{\textwidth}{p{1.4cm} p{4.2cm} p{2.2cm} X}
\toprule
\textbf{Kode} & \textbf{Indikator Kinerja} & \textbf{Target} & \textbf{Metode Pengukuran} \\
\midrule
\textbf{KPI-1} & Akurasi Triase terhadap Dokter & $\mathbf{>92\%}$ & Uji kesesuaian rute triase pada 100 kartu kasus klinis dokter, dengan syarat tidak ada satu pun kondisi darurat yang terlewat. \\
\midrule
\textbf{KPI-2} & Efisiensi Waktu Registrasi Loket & $\mathbf{\ge 40\%}$ & Pemangkasan waktu wawancara awal di loket faskes melalui pemindaian QR rangkuman dibandingkan pendaftaran manual. \\
\midrule
\textbf{KPI-3} & Kepatuhan Puasa Pra-Laboratorium & $\mathbf{>85\%}$ & Proporsi pasien skrining jalur PTM yang datang ke FKTP memenuhi syarat puasa 10--12 jam berkat instruksi pada hasil skrining. \\
\midrule
\textbf{KPI-4} & Kecepatan Deteksi Tanda Bahaya & $\mathbf{<500\text{ ms}}$ & Waktu respon peramban dari ketukan tanda bahaya hingga layar Rute Merah dan tombol 119 tampil di ponsel pasien. \\
\midrule
\textbf{KPI-5} & Skor Kemudahan Pakai (\textit{SUS}) & $\mathbf{>80}$ (\textit{Grade A}) & Evaluasi kuesioner standar \textit{System Usability Scale} pada masyarakat umum dan petugas faskes. \\
\bottomrule
\end{tabularx}
\end{table}


% 6. Bab V: Transparansi Penggunaan AI (Syarat Wajib Hackathon)
%==============================================================
% BAB V TRANSPARANSI PENGGUNAAN KECERDASAN BUATAN
%==============================================================

\section{Transparansi Penggunaan Kecerdasan Buatan}

\subsection{Pernyataan Transparansi AI dalam Siklus Pengembangan}

Sejalan dengan standar integritas akademik dan arahan Kementerian Kesehatan RI mengenai etika kecerdasan buatan, tim pengembang menyampaikan pernyataan terbuka terkait pemanfaatan alat bantu AI generatif dalam proses perancangan, penulisan proposal, dan rekayasa TropiCare. AI diposisikan secara ketat sebagai akselerator kognitif, penyunting tata bahasa, dan asisten kode (\textit{coding copilot}), tanpa menggantikan orisinalitas pemikiran, pertimbangan klinis, maupun tanggung jawab tim manusia (Tabel~\ref{tab:transparansi_ai}).

\begin{table}[H]
\centering
\small
\caption{Matriks Transparansi Pemanfaatan AI dalam Siklus Proyek}\label{tab:transparansi_ai}
\vspace{3pt}
\begin{tabularx}{\textwidth}{p{3.2cm} X X}
\toprule
\textbf{Fase Aktivitas} & \textbf{Kontribusi Alat Bantu AI} & \textbf{Verifikasi \& Pengawasan Manusia} \\
\midrule
\textbf{Perancangan Ide \& Konsep} & Eksplorasi skenario risiko (\textit{edge-cases}), simulasi variasi kosakata awam untuk uji ketahanan leksikal, dan komparasi literatur sistem triase global. & Seleksi kritis oleh tim berdasarkan realitas faskes primer Indonesia dan keselarasan dengan regulasi Kemenkes RI. \\
\midrule
\textbf{Penyusunan Proposal} & Perumusan draf awal kalimat, perbaikan tata bahasa formal, dan penataan struktur matriks komparasi. & Penyuntingan substansial, verifikasi silang dokumen regulasi (UU PDP, UU Kesehatan, KMK PPK FKTP), dan penulisan ulang analisis kausal oleh tim. \\
\midrule
\textbf{Rekayasa Kode \& Desain} & Optimasi sintaksis dokumen LaTeX/TikZ serta generasi kerangka kode dasar (\textit{boilerplate}) FastAPI dan React. & Penulisan mandiri logika inti (algoritma pencocokan KMK, \textit{rule-engine} 11 \textit{Red Flags}, serialisasi QR), tinjauan kode baris-demi-baris, dan pengujian unit. \\
\bottomrule
\end{tabularx}
\end{table}

\subsection{Batasan Sistem dan Pengawasan Manusia (\textit{Human-in-the-Loop})}

Tim TropiCare menegaskan bahwa kecerdasan buatan hanyalah alat bantu (\textit{assistive tool}), bukan pengambil keputusan klinis otonom. Model AI bekerja berbasis probabilitas statistik dan rentan halusinasi jika tidak dibatasi. Oleh karena itu, TropiCare menerapkan tiga pilar pengawasan manusia:

\begin{enumerate}
    \item \textbf{Validasi Medis oleh Kolaborator Klinis:} Seluruh logika pohon keputusan klinis, pembagian 20 penyakit tropis, penetapan 11 kriteria tanda bahaya (\textit{Red Flags}), dan alur rujukan dikurasi langsung oleh kolaborator mahasiswa kedokteran berdasarkan KMK No. HK.01.07/MENKES/1186/2022 \citep{Kemenkes2022a} dan KMK No. HK.01.07/MENKES/1936/2022 \citep{Kemenkes2022b}. Hal ini menjamin kepatuhan pada Standar Kompetensi Dokter Indonesia \citep{KKI2012}.
    \item \textbf{Rekayasa Perangkat Lunak Deterministik:} Arsitektur backend FastAPI, frontend React, basis data direktori 13.930 faskes SATUSEHAT MSI, dan pencarian faskes terdekat dirancang mandiri oleh mahasiswa ilmu komputer tanpa ketergantungan logika otonom pihak ketiga.
    \item \textbf{Keputusan Etis Pembatasan Fitur:} Tim secara sadar menghapus fitur kamera ruam kulit demi mematuhi UU Pelindungan Data Pribadi No. 27/2022 \citep{UUPDP2022} mengenai data biometrik sensitif dan mencegah bias warna kulit sawo matang/gelap, serta menolak penyimpanan nomor KTP/NIK di server demi keamanan privasi masyarakat.
\end{enumerate}

\subsection{Keselarasan dengan Panduan Etika AI Kemenkes RI dan UNESCO}

Operasionalisasi TropiCare berpedoman teguh pada empat pilar etika kecerdasan buatan Kemenkes RI \citep{KemenkesDTO2024} dan rekomendasi UNESCO \citep{UNESCO2021} (Tabel~\ref{tab:etika_ai}).

\begin{table}[H]
\centering
\small
\caption{Keselarasan Implementasi TropiCare dengan Prinsip Etika AI Nasional dan Global}\label{tab:etika_ai}
\vspace{3pt}
\begin{tabularx}{\textwidth}{p{3.5cm} X X}
\toprule
\textbf{Prinsip Etika} & \textbf{Mandat Regulasi} & \textbf{Implementasi Konkret pada TropiCare} \\
\midrule
\textbf{Transparansi \& Keterjelasan (\textit{Explainability})} & AI medis tidak boleh bersifat kotak hitam (\textit{black-box}); penalaran sistem harus dapat diaudit. & Fitur \textbf{``Kenapa aku menebak ini?''} yang menampilkan kutipan langsung kriteria klinis KMK PPK FKTP 2022 sebagai dasar rekomendasi. \\
\midrule
\textbf{Akuntabilitas Mediko-Legal} & Batasan tegas fungsi perangkat lunak dengan wewenang diagnosis dokter berlisensi. & Sistem dibatasi murni pada \textbf{triase pra-hospital non-diagnosis}; sistem tidak pernah meresepkan antibiotik mandiri (mencegah resistensi kuman). \\
\midrule
\textbf{Keamanan Pasien (\textit{Safety-First})} & Pencegahan keterlambatan penanganan akibat kegagalan komputasi. & Deteksi 11 \textit{Red Flags} di peramban yang melompat langsung ke rute darurat IGD/119 dalam $<500\text{ ms}$ tanpa menunggu server. \\
\midrule
\textbf{Pelindungan Privasi (\textit{Privacy-by-Design})} & Penghormatan kedaulatan individu atas data kesehatan pribadi. & Arsitektur \textbf{Zero-NIK}, kalkulasi GPS lokal tanpa disimpan di server, dan pertukaran data klinis via QR Code luring. \\
\bottomrule
\end{tabularx}
\end{table}


% 7. Daftar Pustaka
\section*{DAFTAR PUSTAKA}
\renewcommand{\bibsection}{}
\renewcommand{\bibfont}{\scriptsize}
\setlength{\bibsep}{1.2pt}
\bibliographystyle{ieeetr}
\bibliography{refs/daftar_pustaka}

\end{document}
